\subsection{A conditional spectral enclosure of AUC change over a continuous shrinkage path}
\label{sec:v23-spectral-continuum}
The finite endpoint transition decomposition in
Section~\ref{sec:v22-rank-transitions} can identify rank changes
\emph{after} two scalar-shrinkage choices have been evaluated.
A more demanding question is whether a ranking remains unchanged for
\emph{every} regularisation value on the interval between them.
We exploit the existing source-trained LDA eigensystem, not a
new classifier. For fixed source-trained class contrast $d_s$,
$S=U\operatorname{diag}(s_i)U^\top$ and
$A_\lambda=(1-\lambda)S+\lambda\nu I$,
the score difference on a positive--negative pair of held
\emph{source-validation} trials with vector difference $v_{pn}$ is
\begin{equation}
 z_{pn}(\lambda)=v_{pn}^\top A_\lambda^{-1}d_s
 =\sum_{i=1}^p\frac{(u_i^\top v_{pn})(u_i^\top d_s)}
 {(1-\lambda)s_i+\lambda\nu}.
 \label{eq:v23-spectral-path}
\end{equation}
Every candidate is required to have a positive denominator over the
entire interval. For any subinterval $[l,r]$ and coefficient
$a_{i,pn}=(u_i^\top v_{pn})(u_i^\top d_s)$, each individual rational
summand is monotone because its denominator is affine and strictly
positive. Hence for all $\lambda\in[l,r]$,
\begin{align}
 L_{pn}[l,r]&=\sum_i\min\left\{
 \frac{a_{i,pn}}{(1-l)s_i+l\nu},
 \frac{a_{i,pn}}{(1-r)s_i+r\nu}\right\},\nonumber\\
 U_{pn}[l,r]&=\sum_i\max\left\{
 \frac{a_{i,pn}}{(1-l)s_i+l\nu},
 \frac{a_{i,pn}}{(1-r)s_i+r\nu}\right\},
 \label{eq:v23-interval}
\end{align}
enclose $z_{pn}(\lambda)$ as
$L_{pn}[l,r]\le z_{pn}(\lambda)\le U_{pn}[l,r]$.
Consequently, if every subinterval has $L_{pn}>0$, that pair
remains positive throughout the complete path; similarly if
every subinterval has $U_{pn}<0$, its sign remains negative.
Pairs that fail the test are \emph{uncertain}, not observed
reversals. We additionally intersect the termwise enclosure
with a centered second-order Taylor interval.
For $m=(l+r)/2$, $h=(r-l)/2$, define
$z_m=z_{pn}(m)$, $z'_m$ from the spectral derivative,
and the curvature envelope
\[
 B_{2,pn}[l,r]=
 2\sum_i\frac{|a_{i,pn}|(\nu-s_i)^2}
 {\min\{(1-l)s_i+l\nu,(1-r)s_i+r\nu\}^{3}}.
\]
Taylor's theorem implies
\[
 z_{pn}(\lambda)\in
 [z_m-R_{pn},z_m+R_{pn}],\quad
 R_{pn}=h|z'_m|+\tfrac12h^2B_{2,pn}[l,r].
\]
Intersecting the two valid intervals retains the enclosure
while exploiting first-derivative cancellation. This refinement
was developed \emph{after} the first continuous-search pilot
showed a loose residual gap and is not an independently
prospectively confirmed improvement.

\paragraph{Finite-trial AUC sensitivity bound.}
For a fixed validation person $g$, denote by $\eta_g$
the fraction of positive--negative trial pairs not certified
strictly sign-stable over the whole interval. Since every
certified pair contributes unchanged AUC credit $0$ or $1$
throughout, while each uncertified pair may vary by no more
than one unit of credit, for all $\lambda\in[l,r]$,
\begin{equation}
 \left|A_g(\lambda)-A_g(l)\right|\le\eta_g.
 \label{eq:v23-continuum-bound}
\end{equation}
This is a deterministic bound on a fixed participant's fixed
sample scores, conditional on the computed source eigensystem;
it is \textbf{not} an unseen-person generalisation guarantee
or a statistical confidence interval. The mathematical
ingredients---spectral shrinkage and scalar interval arithmetic---
are standard. We do not claim novelty for either ingredient.
In implementation, a conservative floating-point
slack prevents a near-zero pair from receiving a fragile
rank-stability certificate; without machine-verified directed
rounding we do not claim formal verified-numerics correctness
for the original real-valued dataset.

The procedure is evaluated on source-validation subjects
under the same original fixed EEG covariance and
tangent-space folds as V21--V22. The candidate $\lambda$
endpoints are computed from SOURCE people only; no outer
target data enter this analysis. Crucially, decreasing
the bound in Eq.~\eqref{eq:v23-continuum-bound}
is a property of the source trial-score path, not evidence
that a future target subject will be protected. To promote
this analysis to a novel useful model-selection policy
requires a frozen source-statistic-to-selection map and an
unexposed external acquisition that tests actual target harm.
